\section{Information categories, local periodicity and retained results}
\label{sec:comparison}
\subsection{What is observed and what is reconstructed}
There are three different collision experiments in the present submission
package. The spatial first-impact experiment of Supplement P launches
uniformly in a fixed laboratory square with independent uniform directions
and records a localized first-impact position. In exact data, the support
of its position law is the relevant physical boundary union. Its protected
launch collars give positive mass to every boundary arc. The exact
finite-patch inverse in that experiment is therefore informationally
direct; it is not boundary recovery from hidden travel-time or spectral
data.

The experiment of this primary article records no collision coordinates.
It instead allows localized commanded launch densities and two opposed
collimated directions. The measured imbalance is a finite difference
of occupation, not the occupation function itself. The cancellation in
Theorem~\ref{thm:balance} removes paths which enter and leave a body
between two free endpoints. The finite-chain zero condition restores
the additive information lost in taking that difference. The same
nonlinear inverse works for mollified occupation without differentiating
noisy probability functions. The theorem is thus a scalar-output
alternative within the collision-law programme, with an explicitly
richer preparation contract than a single uniform launch law. Neither
experiment is identified with the other by an unsupported information
ordering.

Earlier intrinsic return-law results concern different selected
itineraries, endpoint arclength laws and moment compressions. Their
count-only finite fibers, formal fibers and smooth Taylor agreements
remain distinct from equality of exact analytic count germs. A family
of spatially commanded binary collision probabilities is not the
previously studied uniformly prepared count-only datum. The present
argument makes no inference about that analytic count-germ problem,
passive trajectories, marked-length spectra or spectral rigidity.

\subsection{Finite differences and local periodicity}
The finite-difference operator in \eqref{eq:balance} is elementary.
The point is the physical identity realizing it without output positions,
and the bounded zero-witness inverse, rather than a new general theorem
about difference operators. Covariogram formulas also link translations
to occupation statistics: for a measurable finite-volume set, the
covariogram integrates the product of two translated indicators.
Galerne~\cite{Galerne} relates its directional derivatives to perimeter
and directional variation. Here spatially specified test densities
retain the location of a \emph{signed difference}; the datum is not a
translation-invariant autocorrelation. No covariogram uniqueness or
priority conclusion follows from the comparison.

There is also a distinction between recognizing a period and proving
that a set must be periodic. Lagarias and Pleasants~\cite{LP} study
patch-counting complexity and crystallinity of Delone sets. Dolbilin,
Garber, Schulte and Senechal~\cite{DGSS} study regularity radii: local
cluster equivalence can impose global regularity, with quantitative
bounds in the classes they treat. Herva and Kari~\cite{HK} develop
local-complexity and annihilator methods for Delone configurations,
including periodic decomposition and criteria forcing periodicity.
These are local-to-global structural questions with their own hypotheses.

In contrast, \eqref{eq:period-presentation} already assumes a bounded
periodic presentation. Its bound ensures a known finite central patch
contains representatives of every orbit and a protected list containing
generators of the full period group. Our period test does not infer
crystallinity for an arbitrary Delone set from observing a large patch.
The new scalar balance need not vanish and is not an annihilator
hypothesis. Its inversion first recovers the physical patch; only then
is the previously proved finite translation criterion used. This
separation places the contribution without suggesting that bounded
period recognition supersedes general local-periodicity theory.

\subsection{Two finite experiments and their different resources}
On the repeated-motif patch-margin class, the retained spatial first-hit
theorem gives a sufficient attempt bound
$C_\eta\nu^{-3/4}\log(C_\eta/(\nu\delta))$, with output localization
$O(\nu^{3/2})$, for its uniform-launch boundary-coverage design. The
scalar theorem here gives
$C_\eta\nu^{-3}\log(C_\eta/(\nu\delta))$, with input localization
$O(\nu^{3/2})$ and no position output, using a two-dimensional grid of
commands and a fixed precision per scalar mean. Both constants and
accuracy thresholds depend on the same kind of nonperiod patch margin,
in addition to their respective apparatus and geometric bounds. These
are sufficient upper bounds for different experiments, not competing
minimax rates. The former's boundary coverage and the latter's scalar
spatial grid count different operations.

The boundary smoothing estimate is a compensated-convolution argument,
and the probability estimate is an elementary bounded-variable inequality.
Neither exponent is offered as a separate assertion of global optimality.
The exact reversal identity, finite-chain inverse and the explicit
preparation/output trade are the additions; the geometric period lemma
and arithmetic classification retain their preceding provenance.

\subsection{Submission architecture and source status}
Supplement P is the entire reviewed A2 v28 paper and its nested source
history, supplied without changing any mathematical input. It includes
the repeated-motif and translation-separated first-impact theorems,
record-local collision clouds, intrinsic return-law inverses, completion,
sequential certificates, moment theory and the relative-law supplementary
volumes. Those results are neither deleted nor asserted anew by the
primary article. The present text includes the period proof it uses
and has no dependency on the older return-law chapters.

The phrase ``already published'' in the preserved v28 README refers to
a repository deposit and must not be read as journal publication of v27.
The corrected status in this submission is \emph{previously deposited
and source-qualified author revision}. A prior exact-source build
qualifies its own frozen source only; the new primary and retained
volumes need the current qualification receipt. These source statements
do not replace independent proof review or a publication record.
